\documentclass{situationist-preprint}

\title{Meaning Without Achievement}
\msid{EXI-17-a}
\msfamily{EXI-17}
\msdiscipline{Meaning, Death and the Absurd}
\mstype{Principal manuscript}
\mskeywords{meaning in life; achievement; engagement; fitting fulfilment; Wolf; failure; experience machine}

\begin{document}
\maketitle

\begin{abstract}
Does a meaningful life require achievement? Many accounts of meaning give success a central place: a meaningful life is one that makes a difference, leaves traces, or accomplishes something difficult and worthwhile. We argue that it does not. A meaningful life requires active engagement with something worthwhile, in the sense of Wolf's fitting-fulfilment view, and achievement is not necessary. The argument turns on pairs of lives that are alike in engagement and in the worth of their objects and differ only in success: the mathematician who spends her career on a problem she never solves and the one who solves it, the nurse whose patients all die and the one whose patients recover. Reflection judges the lives in each pair to be meaningful, and judges the unsuccessful life to be far from meaningless. Achievement adds a further good, and may add to the degree of meaning; it is not a condition of meaning. We distinguish the requirement of contact with reality, which the experience machine shows to matter, from a requirement of success, which it does not, and show that the view secures the universal availability of meaning.
\end{abstract}

\section{The problem}

A widespread thought about meaning is that it is tied to accomplishment. On this thought, the lives we cite as paradigms of meaning, those of Einstein, Mandela or Mother Teresa in the examples that dominate the literature \citep{Metz2013}, are meaningful because of what they brought about, and a life that brings nothing about is to that extent empty. \citet[ch.~6]{Nozick1981} connected meaning with transcending one's limits by making a difference to something beyond oneself, and with leaving traces. \citet{Bradford2015} has given an account of achievement as the competent causation of an outcome through a difficult process, and has argued that achievement is intrinsically valuable; it is natural to suppose that this is the value that meaning tracks.

The opposite view is that meaning is a matter of how one engages with what one does. \citet{Taylor1970} imagined that the gods, in pity, implanted in Sisyphus a compulsive desire to roll stones, and concluded that his life would then have all the meaning a life can have, though he achieved nothing. Taylor's subjectivism was rightly rejected: a life devoted with passion to rolling stones, or to counting blades of grass \citep[\S65]{Rawls1971}, is not meaningful. \citet{Wolf1997,Wolf2010} combined the subjective and objective elements. On her fitting-fulfilment view, meaning arises when subjective attraction meets objective attractiveness: the subject must be actively engaged, and the object of her engagement must be worthwhile independently of her engagement.

Wolf's view does not mention success. Whether it should is our question. We argue that it should not: meaning requires active engagement with something worthwhile, and achievement is not necessary.

\section{Framework}

\begin{definition}[Active engagement]\label{d:engage}
An agent is \emph{actively engaged} with an object if she acts in ways directed at it, cares how it fares, and is gripped or absorbed by it, so that her activity expresses her concern rather than being merely performed.
\end{definition}

\begin{definition}[Worthwhile]\label{d:worth}
An object of engagement (a project, a person, a practice, an inquiry) is \emph{worthwhile} if it has or serves value that does not depend on the agent's attitudes toward it.
\end{definition}

On Definition~\ref{d:worth}, the worthwhile includes the components of the good identified in ETH-01: well-being, knowledge, achievement, relationships, fairness, beauty and autonomy. Achievement is one of these goods. Our thesis is not that it lacks value, but that meaning does not require it.

\begin{definition}[Achievement]\label{d:ach}
An agent \emph{achieves} an outcome if she competently brings it about through a process that is difficult for her \citep{Bradford2015}. An activity is \emph{successful} if it attains the outcome at which it is directed.
\end{definition}

\begin{principle}[Fitting fulfilment]\label{p:ff}
A life is meaningful to the extent that its subject is actively engaged with what is worthwhile.
\end{principle}

\begin{principle}[Contact with reality]\label{p:real}
Active engagement in the sense relevant to Principle~\ref{p:ff} requires that the agent's activity be directed at a real object and actually bear on it; engagement with an object that exists only in the agent's experience does not count.
\end{principle}

Principle~\ref{p:ff} is Wolf's view. Principle~\ref{p:real} applies to meaning the lesson that ETH-05 draws for well-being from Nozick's experience machine: a life of experiences as of curing diseases, lived in the machine, is not a life engaged with the sick. Principle~\ref{p:real} requires that one's activity reach its object; it does not require that it succeed.

\begin{assumption}[Reflective judgment]\label{a:refl}
Considered judgments about whether lives are meaningful, made under idealized reflection with the relevant facts in view, are evidence about meaning, in the same way that such judgments are evidence about the good in ETH-01.
\end{assumption}

\section{The argument}

We consider pairs of lives that are alike in every respect relevant to Principle~\ref{p:ff} except the success of the activity.

\paragraph{The unsolved problem.} Two mathematicians devote their careers to the same open problem, which is important to their field. Each works with absorption, cares about the result, and brings to the work comparable competence and ingenuity. One proves the conjecture. The other does not, and the problem remains open at her death.

\paragraph{The ward.} Two nurses work in different wards with equal skill and devotion. In the first, most patients recover. In the second, a ward for patients with a disease for which there is no treatment, all patients die; the nurse eases their suffering and keeps them company.

\paragraph{The listener.} A person spends much of her life listening to and studying music she loves, with growing understanding, but produces nothing: she does not perform, compose or write about it.

\begin{lemma}[Failure is compatible with meaning]\label{l:fail}
The unsuccessful mathematician's life is meaningful, and so is the life of the nurse in the second ward.
\end{lemma}

\begin{proof}
Each is actively engaged in the sense of Definition~\ref{d:engage}. Each object of engagement is worthwhile in the sense of Definition~\ref{d:worth}: the problem matters to knowledge, which is a component of the good, and the patients' well-being matters whether or not it can be restored. Each activity is directed at a real object and bears on it, so Principle~\ref{p:real} is satisfied: the mathematician's work rules out approaches and clarifies the problem, and the nurse's care changes how her patients live their last weeks. By Principle~\ref{p:ff}, both lives are meaningful. Reflective judgment, which Assumption~\ref{a:refl} counts as evidence, concurs: it is not the verdict of reflection that a life devoted to an unsolved problem, or to the care of the dying, was meaningless.
\end{proof}

\begin{lemma}[Meaning without product]\label{l:listener}
The listener's life is meaningful to the extent of her engagement with music.
\end{lemma}

\begin{proof}
Her engagement is active, since study and attentive listening are activities directed at their object; the object, music of value, is worthwhile; her listening bears on that object, which is real. By Principle~\ref{p:ff}, her engagement contributes meaning. Nothing in her activity is an achievement in the sense of Definition~\ref{d:ach}, since it brings about no outcome beyond her own understanding and pleasure, and the understanding, though hard won, is not the point at which her activity is directed.
\end{proof}

\begin{theorem}\label{t:main}
A meaningful life requires active engagement with something worthwhile; it does not require achievement.
\end{theorem}

\begin{proof}
The first conjunct is Principle~\ref{p:ff}, supported by the failure of the alternatives: subjectivism fails because engagement with worthless objects confers no meaning, and pure objectivism fails because worthwhile work done without engagement, out of mere duty or inertia, is not what reflection counts as a meaningful life \citep{Wolf2010}. For the second conjunct: if achievement were necessary for meaning, the unsuccessful mathematician's life, the life of the nurse in the second ward, and the listener's life would lack it, so far as their central activities are concerned. By Lemmas~\ref{l:fail} and~\ref{l:listener} they do not lack it.
\end{proof}

\begin{corollary}[Achievement is neither necessary nor sufficient]\label{c:neither}
Achievement is not sufficient for meaning either: an achievement in a worthless pursuit, or one produced without engagement, confers none.
\end{corollary}

\begin{proof}
Counting blades of grass faster than anyone has before may be difficult and competently done, and is therefore an achievement by Definition~\ref{d:ach}, but its object is not worthwhile. A worthwhile result obtained by someone wholly indifferent to it fails the condition of engagement. By Principle~\ref{p:ff}, neither confers meaning.
\end{proof}

\begin{remark}[Degree]\label{r:degree}
Theorem~\ref{t:main} is compatible with the claim that, other things equal, the successful mathematician's life is more meaningful than the unsuccessful one's, or at least better. Achievement is a component of the good (ETH-01), and it adds value to a life; if success enlarges the worth realized through one's engagement, it may enlarge meaning too. We take no position on that question. What the pairs show is that the difference between the lives in each pair is not the difference between a meaningful and a meaningless life.
\end{remark}

\begin{corollary}[Availability]\label{c:avail}
Meaning so understood is not hostage to luck in outcomes.
\end{corollary}

\begin{proof}
Whether an activity succeeds depends on circumstances outside the agent's control: on whether the problem is solvable, the disease treatable, the market favorable. Whether she is engaged with something worthwhile depends far less on such circumstances. By Theorem~\ref{t:main}, meaning depends on the second and not on the first.
\end{proof}

Corollary~\ref{c:avail} supports, without by itself establishing, the conclusion of EXI-01 that meaning in this sense is available to everyone.

\section{Objections and replies}

\begin{objection}[Everything is an achievement]\label{o:broad}
On Bradford's account, any difficult and competent activity is an achievement, since the activity itself is an outcome. The mathematician who fails has achieved years of rigorous work; the listener has achieved understanding. So every meaningful life contains achievement after all.
\end{objection}

\begin{reply}
If achievement is understood so broadly that the competent performance of a difficult activity counts, then active engagement in a demanding worthwhile activity entails achievement trivially, and the thesis that meaning requires achievement becomes a restatement of the engagement condition. The substantive thesis we deny is that meaning requires success: attaining the outcome at which one's activity is directed. The unsuccessful mathematician did not attain it, and her life is meaningful. Not all engagement is difficult, moreover: love of a friend need not be.
\end{reply}

\begin{objection}[Making a difference]\label{o:nozick}
Meaning requires connection with value beyond oneself, and the way to connect with something is to make a difference to it \citep{Nozick1981}. A life that changes nothing is enclosed within itself.
\end{objection}

\begin{reply}
We accept that meaning requires connection with value beyond oneself; Principles~\ref{p:ff} and~\ref{p:real} require it. But connection is not success. The nurse in the second ward is connected with her patients as closely as anyone can be, and she makes a difference to how they die, though not to whether they die. The listener is connected with music through attention and understanding. Connection requires that one's activity reach its object, not that it attain its goal.
\end{reply}

\begin{objection}[The experience machine]\label{o:machine}
ETH-05 shows that a life in contact with reality is better than its experiential duplicate in which the subject's aims are not really achieved. This is an argument that success matters. Why does it not show that meaning requires success?
\end{objection}

\begin{reply}
ETH-05 concerns well-being, and its comparison contrasts a veridical life with an experiential duplicate. Two features separate them: veridicality and the achievement of aims. Our pairs separate the two features. The unsuccessful mathematician is not in a machine: her experiences are veridical, her work bears on a real problem, and her life is in contact with reality. Principle~\ref{p:real} carries over to meaning the part of ETH-05's lesson that concerns reality. The part that concerns success bears on how well her life goes for her, and we grant that success would have made it go better; it does not bear on whether her life is meaningful.
\end{reply}

\begin{objection}[Sisyphus]\label{o:sisyphus}
If achievement is unnecessary, \citeauthor{Taylor1970}'s Sisyphus, engaged with his stone and achieving nothing, lives a meaningful life.
\end{objection}

\begin{reply}
He does not, but not because he achieves nothing: his object is not worthwhile. A Sisyphus whose stone rolling kept a village from flooding would live a meaningful life, even if the stone always rolled back and the village was always threatened again. Taylor's case fails the condition of worth, and the conclusion it is used to support is the subjectivism we reject.
\end{reply}

\section{Conclusion}

A meaningful life requires active engagement with something worthwhile. It does not require that the engagement succeed. Pairs of lives that differ only in success, such as those of two mathematicians or two nurses, are judged on reflection to be meaningful in both cases; and lives devoted to appreciation rather than production are meaningful without achieving anything. Achievement is a good, and may enlarge meaning, but it is neither necessary nor sufficient for it. What the experience machine shows is that engagement must reach a real object, not that it must attain its end.

\begin{thebibliography}{9}\small
\bibitem[Bradford(2015)]{Bradford2015} Bradford, G. (2015). \emph{Achievement}. Oxford University Press.
\bibitem[Metz(2013)]{Metz2013} Metz, T. (2013). \emph{Meaning in Life: An Analytic Study}. Oxford University Press.
\bibitem[Nozick(1981)]{Nozick1981} Nozick, R. (1981). \emph{Philosophical Explanations}. Harvard University Press.
\bibitem[Rawls(1971)]{Rawls1971} Rawls, J. (1971). \emph{A Theory of Justice}. Harvard University Press.
\bibitem[Taylor(1970)]{Taylor1970} Taylor, R. (1970). \emph{Good and Evil: A New Direction}. Macmillan.
\bibitem[Wolf(1997)]{Wolf1997} Wolf, S. (1997). Happiness and meaning: Two aspects of the good life. \emph{Social Philosophy and Policy}, 14(1), 207--225.
\bibitem[Wolf(2010)]{Wolf2010} Wolf, S. (2010). \emph{Meaning in Life and Why It Matters}. Princeton University Press.
\end{thebibliography}

\end{document}
